\chapter{Appendix 附錄}
\label{ch:appendix}

本附錄整理撰寫論文時常見的排版與語法注意事項，供大家撰寫與校稿時參考，請記得看latex的語法。(此段落主要由ChatGPT撰寫整理 XD)

\section*{文字與標點}
\begin{itemize}
  \item 英文引號請使用成對反引號與直引號：``text'' → “text”。
  \item 破折號（dash）有三種：`-`（連字號，如 \textit{pre-trained}）、`--`（範圍，如 10--20）、`---`（句中強調，如 the result---as expected---was correct）。
  \item 省略號請使用 \verb|\ldots| 而非三個句點：A, B, C\ldots, Z。
  \item 數值與單位間需留不換行空白：10~km，30$^\circ$C。
  \item 中英文混排時，中文與英文之間宜留一半形空白，例如：中文 English 中文。
\end{itemize}

\section*{數學與符號}
\begin{itemize}
  \item 數學變數自動為斜體，單位或文字常數須使用 \verb|\mathrm{}| 或 \verb|\text{}|，如 $v=10~\mathrm{m/s}$。
  \item 三角函數、極值運算等應以數學命令表示：$\sin\theta$, $\max(x,y)$。
  \item 向量或矩陣可用粗體：$\mathbf{x}=[x_1,x_2]^\top$。
  \item 多行方程式請使用 \texttt{align} 環境，勿以多個 \texttt{equation} 拆寫。
\end{itemize}


\section*{圖表與引用}
\begin{itemize}
  \item 圖片標題（caption）置於圖下方
  \item 表格標題置於表上方。
  \item 參照圖表時請使用 \verb|\ref{}| 或 \verb|\autoref{}|，勿手動輸入編號。\\
        例：As shown in Fig.~\ref{fig:excel2latex}.
  \item 表格建議使用 \texttt{booktabs} 套件：\verb|\toprule|、\verb|\midrule|、\verb|\bottomrule|。
\end{itemize}

\section*{文獻與引用}
\begin{itemize}
  \item 文獻引用使用 \verb|\cite{key}| 或 \verb|\parencite{key}|。
  \item 多篇同時引用可寫為 \verb|\cite{a,b,c}| → [1–3]。
  \item 特殊大小寫（如 GPU, LaTeX）須以大括號保留：\verb|{GPU}|。
\end{itemize}

\section*{Related Work Reference List}
\begin{itemize}
  \item[{[1]}] Cox, S. J., \& Dobson, D. C. (1999). Maximizing band gaps in two-dimensional photonic crystals. \textit{SIAM Journal on Applied Mathematics}, 59(6), 2108-2120.
  \item[{[2]}] Borel, P. I., Harpøth, A., Frandsen, L. H., Kristensen, M., Shi, P., Jensen, J. S., \& Sigmund, O. (2004). Topology optimization and fabrication of photonic crystal structures. \textit{Optics Express}, 12(9), 1996-2001.
  \item[{[3]}] Jensen, J. S., \& Sigmund, O. (2011). Topology optimization for nano-photonics. \textit{Laser \& Photonics Reviews}, 5(2), 308-321.
  \item[{[4]}] Lalau-Keraly, C. M., Bhargava, S., Miller, O. D., \& Yablonovitch, E. (2013). Adjoint shape optimization applied to electromagnetic design. \textit{Optics Express}, 21(18), 21693-21701.
  \item[{[5]}] Piggott, A. Y., Lu, J., Lagoudakis, K. G., Petykiewicz, J., Babinec, T. M., \& Vu\v{c}kovi\'{c}, J. (2015). Inverse design and demonstration of a compact and broadband on-chip wavelength demultiplexer. \textit{Nature Photonics}, 9(6), 374-378.
  \item[{[6]}] Shen, B., Wang, P., Polson, R., \& Menon, R. (2015). An integrated-nanophotonics polarization beamsplitter with $2.4\times 2.4\,\mu\text{m}^2$ footprint. \textit{Nature Photonics}, 9(6), 378-382.
  \item[{[7]}] Frellsen, L. F., Ding, Y., Sigmund, O., \& Frandsen, L. H. (2016). Topology optimized mode multiplexing in silicon-on-insulator photonic wire waveguides. \textit{Optics Express}, 24(15), 16866-16873.
  \item[{[8]}] Piggott, A. Y., Petykiewicz, J., Su, L., \& Vu\v{c}kovi\'{c}, J. (2017). Fabrication-constrained nanophotonic inverse design. \textit{Scientific Reports}, 7(1), 1786.
  \item[{[9]}] Molesky, S., Lin, Z., Piggott, A. Y., Jin, W., Vu\v{c}kovi\'{c}, J., \& Rodriguez, A. W. (2018). Inverse design in nanophotonics. \textit{Nature Photonics}, 12(11), 659-670.
  \item[{[10]}] Hammond, A. M., Oskooi, A., Johnson, S. G., \& Ralph, S. E. (2021). Photonic topology optimization with semiconductor-foundry design-rule constraints. \textit{Optics Express}, 29(15), 23916-23938.
  \item[{[11]}] Pao, Yu-Cheng; Wu, Jhih-Sheng. Nano-Focusing by Topology Optimization. (2026).
  \item[{[12]}] Kirkpatrick, S., Gelatt, C. D., \& Vecchi, M. P. (1983). Optimization by simulated annealing. \textit{Science}, 220(4598), 671-680.
  \item[{[13]}] Kadowaki, T., \& Nishimori, H. (1998). Quantum annealing in the transverse Ising model. \textit{Physical Review E}, 58(5), 5355.
  \item[{[14]}] Rendle, S. (2010). Factorization machines. \textit{2010 IEEE International Conference on Data Mining (ICDM)}, 995-1000.
  \item[{[15]}] Kitai, K., Guo, J., Ju, S., Tanaka, S., Tsuda, K., Shiomi, J., \& Tamura, R. (2020). Designing metamaterials with quantum annealing and factorization machines. \textit{Physical Review Research}, 2(1), 013319.
  \item[{[16]}] Inoue, T., Seki, Y., Tanaka, S., Togawa, N., Ishizaki, K., \& Noda, S. (2022). Towards optimization of photonic-crystal surface-emitting lasers via quantum annealing. \textit{Optics Express}, 30(24), 43503-43512.
  \item[{[17]}] Zhu, C., Bamidele, E. A., Shen, X., Zhu, G., \& Li, B. (2024). Machine Learning Aided Design and Optimization of Thermal Metamaterials. \textit{Chemical Reviews}, 124(7), 4258-4331.
  \item[{[18]}] Hu, Jia-Jun; Wu, Jhih-Sheng. Optical Filter and Other Applications Designed by Factorization Machines with Quantum Annealing. (2025).
  \item[{[19]}] Chang, Yi-Hsin; Wu, Jhih-Sheng. Quantum-Compatible Inverse Design of Superoscillatory Lenses. (2026).
  \item[{[20]}] Lazarov, B. S., \& Sigmund, O. (2011). Filters in topology optimization based on Helmholtz-type differential equations. \textit{International Journal for Numerical Methods in Engineering}, 86(6), 765-781.
  \item[{[21]}] Odena, A., Dumoulin, V., \& Olah, C. (2016). Deconvolution and checkerboard artifacts. \textit{Distill}, 1(10), e3.
  \item[{[22]}] Shi, W., Caballero, J., Husz\'{a}r, F., Totz, J., Aitken, A. P., Bishop, R., \ldots \& Wang, Z. (2016). Real-time single image and video super-resolution using an efficient sub-pixel convolutional neural network. \textit{Proceedings of the IEEE conference on computer vision and pattern recognition (CVPR)}, 1874-1883.
  \item[{[23]}] Higgins, I., et al. (2017). $\beta$-VAE: Learning basic visual concepts with a constrained variational framework. \textit{International Conference on Learning Representations (ICLR)}.
  \item[{[24]}] Piggott, A. Y., Petykiewicz, J., Su, L., \& Vu\v{c}kovi\'{c}, J. (2017). Fabrication-constrained nanophotonic inverse design. \textit{Scientific Reports}, 7(1), 1786.
  \item[{[25]}] Burgess, C. P., et al. (2018). Understanding disentangling in $\beta$-VAE. \textit{arXiv preprint arXiv:1804.03599}.
  \item[{[26]}] Ma, W., Cheng, F., Xu, Y., Wen, Q., \& Liu, Y. (2019). Probabilistic representation and inverse design of metamaterials based on a deep generative model with semi-supervised learning strategy. \textit{Advanced Materials}, 31(35), 1901111.
  \item[{[27]}] Zhang, R. (2019). Making convolutional networks shift-invariant again. \textit{International Conference on Machine Learning (ICML)}, 7324-7334.
  \item[{[28]}] Li, Z., Kovachki, N., Azizzadenesheli, K., Liu, B., Bhattacharya, K., Stuart, A., \& Anandkumar, A. (2020). Fourier neural operator for parametric partial differential equations. \textit{ICLR}.
  \item[{[29]}] Kudyshev, Z. A., Kildishev, A. V., Shalaev, V. M., \& Boltasseva, A. (2021). Machine-learning-assisted global optimization of photonic devices. \textit{Nanophotonics}, 10(1), 371-383.
  \item[{[30]}] Karras, T., Aittala, M., Laine, S., H\"{a}rk\"{o}nen, E., Hellsten, J., Lehtinen, J., \& Aila, T. (2021). Alias-free generative adversarial networks. \textit{Advances in Neural Information Processing Systems (NeurIPS)}, 34, 852-863.
\end{itemize}
